\documentclass[11pt]{article}

\usepackage[margin=1in]{geometry}
\usepackage{amsmath, amssymb}
\usepackage{booktabs}
\usepackage{array}
\usepackage{graphicx}
\usepackage{hyperref}
\usepackage{xcolor}
\usepackage{listings}

\lstset{
    basicstyle=\ttfamily\small,
    breaklines=true,
    columns=fullflexible,
    backgroundcolor=\color{gray!8},
    frame=single,
    framerule=0.3pt,
    xleftmargin=2pt,
    xrightmargin=2pt
}

\title{Multi-Robot Double-Angle Formation Controller\\
\large Mathematical Formulation and Data Logging Reference}
\author{TurtleBot3 Swarm Experiments}
\date{\today}

\begin{document}
\maketitle

\section{Overview}

This document describes the controller implemented in
\texttt{multi\_unicycle\_controller\_double\_angle.py}, which drives a group of
$N = 4$ TurtleBot3 unicycle robots into and around a desired formation while
simultaneously enforcing two geometric constraints on every controlled
inter-robot pair:

\begin{itemize}
    \item \textbf{Connectivity}: neighbors must not drift further apart than a
    connectivity radius $R_c$ (otherwise they lose line-of-sight / sensing
    range).
    \item \textbf{Safety}: neighbors must not come closer together than a
    safety radius $R_s$ (otherwise their physical bodies risk collision).
\end{itemize}

These constraints are enforced using a \emph{double-angle visual-cone}
parameterization of inter-agent bearing angles, which converts the distance
constraints $R_s + r \le \lVert p_i - p_j \rVert \le R_c - r$ into equivalent
bounds on an angle $\psi_{ij}$ that remains well-conditioned (no singularities)
across the full admissible range.

This document also specifies exactly what data the controller now logs to
CSV (replacing the previous GIF/animation export, which was slow and
memory-heavy for long runs) and what each column means.

\section{Unicycle Model}

Each robot $i$ is modeled as a unicycle with state $(x_i, y_i, \theta_i)$ in
the global frame and control inputs $(v_i, \omega_i)$:

\begin{align}
    \dot{x}_i &= v_i \cos\theta_i \\
    \dot{y}_i &= v_i \sin\theta_i \\
    \dot{\theta}_i &= \omega_i
\end{align}

Odometry feedback provides $(x_i, y_i, \theta_i)$ directly, and body-frame
velocities $(v_i, \omega_i)$ are rotated into global-frame velocities
$(v_{x,i}, v_{y,i}, \omega_i)$ for logging purposes.

\section{Formation Graph and Desired Geometry}

The $4$ robots are connected by $5$ edges, encoded in the incidence matrix
$H \in \mathbb{R}^{5 \times 4}$:

\begin{equation}
    H =
    \begin{pmatrix}
        -1 & 1 & 0 & 0 \\
        0 & -1 & 1 & 0 \\
        0 & 0 & -1 & 1 \\
        1 & 0 & 0 & -1 \\
        0 & 1 & 0 & -1
    \end{pmatrix},
    \qquad
    \bar{H} = H \otimes I_2 \in \mathbb{R}^{10 \times 8}
\end{equation}

Row $k$ of $H$ defines edge $k$ between the two robots whose entries are
$-1$ and $+1$. With $p = [x_1\,y_1\,x_2\,y_2\,x_3\,y_3\,x_4\,y_4]^\top$ the
stacked global positions, the edge (relative-position) vector for all edges is

\begin{equation}
    e = \bar{H}\, p \in \mathbb{R}^{10}, \qquad
    e_k = p_j - p_i \ \text{for edge } k = (i,j).
\end{equation}

A desired formation shape $p^\star$ (see \texttt{self.pstar} in code) defines
the desired edge vectors $e^\star = \bar{H} p^\star$, and the desired
Euclidean distance for edge $k$ is

\begin{equation}
    d_k^\star = \lVert e_k^\star \rVert.
\end{equation}

A second incidence matrix $H_M$ (with $\bar{H}_M = H_M \otimes I_2$) defines
an additional pair used purely for collision-avoidance monitoring between
non-adjacent robots.

\section{Double-Angle Visual-Cone Parameterization}

\subsection{Physical radii and their relationship to the bearing angle}

Let $r$ be the physical radius of each robot footprint, $R_s$ the safety
margin, $R_c$ the connectivity (sensing) radius, and $R_{crit}$ the critical
(hard) collision radius. For a pair of robots whose \emph{centers} are a
distance $d_{ij} = \lVert e_{ij}\rVert$ apart, define the half-angle subtended
by one robot's footprint as seen from the other:

\begin{equation}
    \psi_{ij} = 2 \arcsin\!\left(\frac{r}{d_{ij}}\right).
\end{equation}

This angle increases monotonically as $d_{ij}$ decreases, so it is used as a
proxy for distance that behaves well (no division blow-up) as robots get
close. Three reference angles bound the admissible range of $\psi_{ij}$,
each obtained from the corresponding critical distance via the double-angle
identity $\cos\psi = 1 - 2\sin^2(\psi/2) = 1-2(r/d)^2$:

\begin{align}
    \bar{R}_c &= 1 - 2\left(\frac{r}{R_c - r}\right)^2,
    &\psi_{\min} &= \arccos(\bar{R}_c) \\
    \bar{R}_s &= 1 - 2\left(\frac{r}{R_s + r}\right)^2,
    &\psi_{\max} &= \arccos(\bar{R}_s) \\
    \bar{R}_{crit} &= 1 - 2\left(\frac{r}{R_{crit}}\right)^2,
    &\psi_{crit} &= \arccos(\bar{R}_{crit})
\end{align}

\subsection{Distance bounds used for logging}

Every logged inter-robot distance $d_{ij}$ is compared against the same pair
of physical bounds used to derive the angle limits above:

\begin{equation}
    \boxed{\;\text{lower bound} \;=\; R_s + r\;} \qquad
    \boxed{\;\text{upper bound} \;=\; R_c - r\;}
\end{equation}

\textbf{Why these bounds:} $R_s$ and $R_c$ are specified as clearances/ranges
measured from each robot's \emph{surface}, not its center. Since each robot
occupies a disk of radius $r$, the minimum center-to-center distance that
still preserves a safety clearance of $R_s$ between the two disks' surfaces is
$R_s + r$, and the maximum center-to-center distance that still keeps the two
robots' surfaces within a connectivity range of $R_c$ is $R_c - r$. In the
current configuration ($r = 0.2\,\text{m}$, $R_s = 0.4\,\text{m}$,
$R_c = 3.0\,\text{m}$), this gives:

\begin{equation}
    \text{lower bound} = 0.4 + 0.2 = 0.6\,\text{m}, \qquad
    \text{upper bound} = 3.0 - 0.2 = 2.8\,\text{m}.
\end{equation}

Every edge's desired distance $d_k^\star$ (derived from the desired formation
$p^\star$) is chosen to lie strictly within this interval, and the controller
is designed so that, starting from a feasible configuration, the actual
distance $d_k(t)$ remains within $(\text{lower bound}, \text{upper bound})$
for all $t$.

\subsection{Control law}

For each formation edge $k$ with bearing-derived quantities
$\psi_k$, unit vector $g_k = e_k / d_k$, and left/right rotations of $g_k$ by
$\pm\psi_k/2$:

\begin{equation}
    R(\pm\psi_k/2) =
    \begin{pmatrix}
        \cos(\psi_k/2) & \mp\sin(\psi_k/2) \\
        \pm\sin(\psi_k/2) & \cos(\psi_k/2)
    \end{pmatrix},
    \qquad
    g_{k}^{R} = R(-\psi_k/2)\,g_k,\quad g_{k}^{L} = R(+\psi_k/2)\,g_k
\end{equation}

a projected gradient direction is built from the projector
$P_g = I - g g^\top$:

\begin{equation}
    f_{g,k} = P_{g_k^{L}}\, g_k^{R} + P_{g_k^{R}}\, g_k^{L}, \qquad
    f_{\psi,k} = \frac{\sin(2\psi_k)}{2(1+\cos\psi_k)}.
\end{equation}

A barrier-like gain $\xi_k$ pushes the edge back toward $d_k^\star$ as
$\psi_k$ approaches either limit angle, using the error term
$E_k = \cos^2\psi_k - \cos^2\psi_k^\star$ and the corresponding one-sided gain
($K_c$ near the connectivity limit, $K_s$ near the safety limit):

\begin{equation}
    \xi_k = \frac{4 E_k}{K_k^2 - E_k^2}, \qquad
    K_k = q\,K_{c,k} + (1-q)\,K_{s,k}, \qquad
    q = \mathbb{1}[E_k > 0].
\end{equation}

The aggregate formation velocity in the stacked position space is

\begin{equation}
    \dot{p} = -\bar{H}^\top R_\psi\, \xi,
\end{equation}

where $R_\psi = \mathrm{blockdiag}(f_{\psi,1} f_{g,1}, \dots)$. This is then
converted, per robot $i$, into unicycle commands using the current heading
$\theta_i$ and a small look-ahead length $\ell$:

\begin{align}
    v_i &= 2\left(\dot{p}_{2i}\cos\theta_i + \dot{p}_{2i+1}\sin\theta_i\right) \\
    \omega_i &= \frac{-\dot{p}_{2i}\sin\theta_i + \dot{p}_{2i+1}\cos\theta_i}{\ell}
\end{align}

clipped to $v_i \in [-v_{max}, v_{max}]$ and $\omega_i \in [-\omega_{max},
\omega_{max}]$. An analogous but separate term $\xi_\beta$ (using
$R_{crit}$) is computed for pairwise collision avoidance and can optionally be
added to $\dot p$ (disabled by default in the current code).

\section{Data Logging}

\subsection{What changed and why}

The controller previously rendered a combined trajectory + distance
\texttt{matplotlib} animation and saved it as a GIF
(\texttt{animate\_combined}). For long experiments this was slow to render
and consumed large amounts of RAM (every frame is a rasterized image held in
memory before encoding). This has been replaced with a lightweight CSV
export (\texttt{save\_csv}) that streams the already-collected numeric log
into a single file with negligible memory and CPU cost. The PNG summary
plots (final trajectories, distance-vs-time) are still produced, since those
are cheap.

\subsection{Sampling}

State is logged every \texttt{log\_interval} = $0.1\,\text{s}$ of wall-clock
time (see \texttt{log\_trajectories}), once all four robots have produced at
least one odometry message.

\subsection{CSV schema}

The exported file (\texttt{log\_data.csv}, written into the timestamped
output folder alongside the PNG plots) has one row per logged tick and the
following columns:

\begin{center}
\begin{tabular}{@{}l l@{}}
\toprule
\textbf{Column(s)} & \textbf{Description} \\
\midrule
\texttt{time} & Elapsed time since logging started (s) \\
\texttt{bot\_i\_x}, \texttt{bot\_i\_y}, \texttt{bot\_i\_theta} &
    Global-frame pose of robot $i$ for $i \in \{0,1,2,3\}$ (m, m, rad) \\
\texttt{d1} \dots \texttt{d5} &
    Actual inter-robot distance for formation edges $1$--$5$ (m) \\
\texttt{e1} \dots \texttt{e5} &
    Distance error $e_k = d_k - d_k^\star$ for edges $1$--$5$ (m) \\
\texttt{lower\_bound} &
    Constant column, $R_s + r$ (m) --- minimum admissible distance \\
\texttt{upper\_bound} &
    Constant column, $R_c - r$ (m) --- maximum admissible distance \\
\bottomrule
\end{tabular}
\end{center}

Edges $1$--$5$ correspond to rows $1$--$5$ of $H$, i.e.\ the robot pairs
(bot\_0--bot\_1), (bot\_1--bot\_2), (bot\_2--bot\_3), (bot\_3--bot\_0),
(bot\_1--bot\_3) respectively (see \texttt{get\_connected\_pairs}).

\texttt{lower\_bound} and \texttt{upper\_bound} are repeated on every row
(rather than stored once) so that the CSV can be loaded directly into a
plotting tool (e.g.\ \texttt{pandas} + \texttt{matplotlib}) and the bounds
drawn as reference lines against \texttt{d1}\dots\texttt{d5} or
\texttt{e1}\dots\texttt{e5} without any extra bookkeeping.

\subsection{Output directory layout}

On shutdown (\texttt{save\_all\_plots}), the controller writes to
\texttt{\textasciitilde/turtlebot3\_simulations/plots/\allowbreak
<timestamp>/}:

\begin{lstlisting}
~/turtlebot3_simulations/plots/<YYYY-MM-DD_HH-MM-SS>/
    trajectory.png     # final trajectories + initial/final connections
    distances.png      # d_k(t) vs. desired distance, per edge
    log_data.csv        # full numeric log (see schema above)
\end{lstlisting}

\subsection{Post-processing example}

\begin{lstlisting}[language=Python]
import pandas as pd
import matplotlib.pyplot as plt

df = pd.read_csv("log_data.csv")

fig, ax = plt.subplots()
for k in range(1, 6):
    ax.plot(df["time"], df[f"d{k}"], label=f"d{k}")
ax.axhline(df["lower_bound"].iloc[0], color="red", linestyle="--", label="lower bound (Rs+r)")
ax.axhline(df["upper_bound"].iloc[0], color="gray", linestyle="--", label="upper bound (Rc-r)")
ax.set_xlabel("time (s)"); ax.set_ylabel("distance (m)"); ax.legend()
plt.show()
\end{lstlisting}

\end{document}
